% Kompilacija: tectonic reports/izvjestaj.tex   (ili xelatex / lualatex)
\documentclass[11pt]{article}

\usepackage{fontspec}                 % UTF-8 / XeTeX-LuaTeX font handling
\usepackage{polyglossia}
\setdefaultlanguage{croatian}

\usepackage{geometry}
\geometry{a4paper, margin=2.4cm}
\usepackage{graphicx}
\graphicspath{{../outputs/figures/}}
\usepackage{booktabs}
\usepackage{amsmath}
\usepackage{caption}
\captionsetup{font=small, labelfont=bf}
\usepackage{microtype}
\usepackage[hidelinks]{hyperref}
\usepackage{float}

\setlength{\parskip}{0.4em}
\setlength{\parindent}{0pt}

\title{\textbf{Skriveni stilski rizik u portfeljima minimalne varijance}\\[0.3em]
\large Dijagnoza faktorskim klasteriranjem na indeksu Russell 1000 (2005.--2025.)}
\author{Marin Meštrović\\[0.2em]
\small Dubinska analiza podataka i otkrivanje znanja u podatcima --- PMF, Sveučilište u Zagrebu}
\date{Lipanj 2026.}

\begin{document}
\maketitle

\begin{abstract}
\noindent
Portfelji minimalne varijance isključivo s dugim pozicijama uobičajeno se promatraju kao
„neutralna” referenca: po konstrukciji minimiziraju varijancu bez izričitoga stava o
prinosima. Ovaj projekt pokazuje da takvi portfelji na indeksu Russell 1000 nose
\emph{skriven i postojan} nagib prema Fama--French stilskim faktorima profitabilnosti
(RMW) i ulaganja (CMA), te da taj nagib ostaje skriven i sektorskoj i korelacijskoj
diverzifikaciji. Dijagnozu gradimo klasteriranjem dionica u prostoru faktorske
izloženosti --- metodom otkrivanja znanja koja se obrađuje na kolegiju. Klasteriranje
ispravno \emph{lokalizira} problem, no intuitivni ispravak --- ograničenje težine po
klasteru --- pokazuje se neuspješnim. Djelotvoran je ispravak izravno linearno
ograničenje na faktorske bete portfelja, koje ostvarenu koncentraciju stila smanjuje za
31\,\% uz zanemariv trošak u volatilnosti ($p \approx 0.002$).
\end{abstract}

\section{Ciljevi}

Hipoteza projekta jest da mehanička minimizacija rizika neprimjetno daje prednost
određenom području faktorskoga prostora te tako gradi skrivenu stilsku izloženost.
Postavljamo tri pitanja koja ujedno čine i trodijelnu strukturu rada:

\begin{enumerate}
  \item \textbf{Otkriće.} Nose li portfelji minimalne varijance skriven stilski nagib i
        može li se on učiniti vidljivim klasteriranjem u prostoru faktora?
  \item \textbf{Neuspjeli ispravak.} Ako klasteri ocrtavaju problem, ublažava li ga
        ograničavanje težine portfelja po klasteru?
  \item \textbf{Djelotvoran ispravak.} Ako ne, koje ograničenje uistinu smanjuje stilski
        nagib bez žrtvovanja niske volatilnosti?
\end{enumerate}

Metodološko je težište na klasteriranju kao alatu za \emph{otkrivanje znanja}: ono nije
tek opisni korak, nego okvir kroz koji tumačimo cijelu analizu i izvor hipoteze o
ispravku. Uz glavni tijek analize, projekt sustavno obuhvaća metode kolegija ---
particijsko, vjerojatnosno i na gustoći utemeljeno klasteriranje, odabir broja klastera
indeksima interne validacije te nadzirana stabla odlučivanja i ansamble.

\section{Podatci}

\begin{itemize}
  \item \textbf{Univerzum.} Trenutačni sastav indeksa Russell 1000 (1.002 dionice s
        pripadajućim GICS sektorima). Pristranost preživljavanja priznajemo: rabi se
        trenutačno članstvo u indeksu.
  \item \textbf{Cijene.} Dnevne prilagođene zaključne cijene (Yahoo Finance, \texttt{yfinance}),
        agregirane na zadnji dan u mjesecu.
  \item \textbf{Faktori.} Fama--French petfaktorski mjesečni prinosi iz zbirke podataka
        Kennetha Frencha (MKT, SMB, HML, RMW, CMA).
  \item \textbf{Razdoblje.} Od siječnja 2005. do prosinca 2025.\ (21 godina). Prvih pet
        godina čini prvi prozor procjene, a testne godine teku od 2010.\ do 2025.
\end{itemize}

Univerzum raste s približno 600 dionica u prozoru 2009-12 na približno 880 u prozoru
2024-12, kako dionice s kasnijim IPO-om počinju ispunjavati uvjete (slika~\ref{fig:coverage}).

\begin{figure}[H]
  \centering
  \includegraphics[width=0.78\linewidth]{01_universe_coverage.png}
  \caption{Pokrivenost univerzuma kroz klizne prozore procjene: broj dionica s najmanje
  60 valjanih mjeseci i čistom petfaktorskom prilagodbom.}
  \label{fig:coverage}
\end{figure}

\section{Metodologija}

\paragraph{Značajke dionice.} Za svaki par \emph{(prozor, dionica)} procjenjujemo OLS
petfaktorsku regresiju (metoda najmanjih kvadrata) na viškovima prinosa i bilježimo
$(\alpha,\ \beta_{\text{MKT}},\ \beta_{\text{SMB}},\ \beta_{\text{HML}},\
\beta_{\text{RMW}},\ \beta_{\text{CMA}},\ \sigma_\varepsilon,\ R^2)$. Prilagodbe koje krše
granice kvalitete ($|\beta|\le 5$, $\sigma_\varepsilon\le 0.50$, $R^2\ge 0.05$) izbacuju se
iz toga prozora.

\paragraph{Klasteriranje.} Glavni predmet istraživanja jesu \emph{faktorski klasteri}:
Wardova veza na euklidskoj udaljenosti nad šest standardiziranih značajki (z-vrijednosti)
$(\beta_{\text{MKT}}, \beta_{\text{SMB}}, \beta_{\text{HML}}, \beta_{\text{RMW}},
\beta_{\text{CMA}}, \sigma_\varepsilon)$. Kao referentne podjele za usporedbu služe
\emph{korelacijski klasteri} (potpuna veza na udaljenosti $\sqrt{2(1-\rho)}$) i
\emph{GICS sektori}. Broj klastera $K=7$ najveća je vrijednost po prosječnoj silueti koja
zadovoljava dva uvjeta uloživosti: svaki klaster ima $\ge 4$ dionice u svakom prozoru te
$K \times \text{group\_cap} \ge 1$ uz $\text{group\_cap}=0.15$.

\paragraph{Portfelji.} Svi su portfelji isključivo s dugim pozicijama, u svakom prozoru
potpuno uloženi, uz ograničenje pojedinačne pozicije $w_{\max}=0.02$:

\begin{enumerate}
  \item jednake težine $1/N$;
  \item minimalna varijanca bez grupnoga ograničenja;
  \item minimalna varijanca uz $\sum w_i \le 0.15$ unutar svakog GICS sektora;
  \item isto grupno ograničenje na korelacijskim klasterima;
  \item isto grupno ograničenje na faktorskim klasterima;
  \item faktorski neutralna minimalna varijanca:
        $|w^{\!\top}\beta_f|\le\varepsilon$ za $f\in\{\text{SMB, HML, RMW, CMA}\}$,
        uz $\varepsilon\in\{0.00, 0.05, 0.10, 0.15\}$.
\end{enumerate}

Kovarijancu procjenjujemo Ledoit--Wolfovim sažimanjem na mjesečnim prinosima iz prozora
procjene, preračunatu na godišnju razinu množenjem s 12. Glavna mjera nagiba jest
\textbf{koncentracija stila}
$\;|\beta_{\text{SMB}}| + |\beta_{\text{HML}}| + |\beta_{\text{RMW}}| + |\beta_{\text{CMA}}|$.

\paragraph{Klizni unaprijedni postupak (\emph{walk-forward}) i bootstrap.} Prozor procjene
obuhvaća 60 mjeseci unatrag, testni horizont 12 mjeseci, a ponovnu procjenu provodimo
svakih 12 mjeseci $\rightarrow$ 16 disjunktnih testnih godina (2010.--2025.). Za intervale
pouzdanosti razlika među metrikama koristimo blok-bootstrap (pomični blokovi od 12
mjeseci, 1000 ponovnih uzoraka).

\section{Čin I --- Otkriće}

\subsection{Klasteriranje otkriva faktorsku strukturu univerzuma}

Faktorsko klasteriranje dijeli univerzum na sedam stilski koherentnih skupina.
Radar-grafovi centroida (slika~\ref{fig:radar}) prikazuju njihove standardizirane faktorske
profile, a UMAP-projekcija, obojena na tri načina (slika~\ref{fig:umap}), otkriva ključan
nalaz: faktorski klasteri dijele univerzum po granicama koje ne prepoznaju ni korelacijski
klasteri ni GICS sektori.

\begin{figure}[H]
  \centering
  \includegraphics[width=0.92\linewidth]{02_cluster_centroids_radar.png}
  \caption{Radar-grafovi centroida sedam faktorskih klastera (standardizirani faktorski
  profili). Svaki klaster odgovara prepoznatljivu stilu.}
  \label{fig:radar}
\end{figure}

\begin{figure}[H]
  \centering
  \includegraphics[width=\linewidth]{02_factor_umap_three_panels.png}
  \caption{UMAP-projekcija prostora faktora, obojena po faktorskom klasteru, korelacijskom
  klasteru i GICS sektoru. Faktorska se geometrija ne podudara ni sa sektorima ni s
  korelacijskom strukturom.}
  \label{fig:umap}
\end{figure}

Stabilnost klastera pri bootstrap-ponovnom uzorkovanju umjerena je: srednji Jaccardov
indeks zajedničke pripadnosti iznosi $0.30$ za faktorske naspram $0.26$ za korelacijske
klastere, dovoljno da potvrdi kako struktura nije slučajnost svojstvena uzorku.

\subsection{Skriveni nagib: minimalna varijanca naginje profitabilnosti i niskom ulaganju}

Petfaktorska atribucija prinosa na punom testnom uzorku otkriva nagib (tablica~\ref{tab:attr1}).

\begin{table}[H]
  \centering
  \caption{Petfaktorska atribucija na punom uzorku (192 mjeseca).}
  \label{tab:attr1}
  \small
  \begin{tabular}{lrrrrrr}
    \toprule
    Portfelj & $\beta_{\text{MKT}}$ & $\beta_{\text{SMB}}$ & $\beta_{\text{HML}}$ &
    $\beta_{\text{RMW}}$ & $\beta_{\text{CMA}}$ & konc.\ stila \\
    \midrule
    Jednake težine                 & 1.11 & 0.49 & +0.09 & 0.08 & 0.17 & 0.83 \\
    Minimalna varijanca            & 0.66 & 0.20 & $-0.12$ & \textbf{+0.33} & \textbf{+0.34} & 0.98 \\
    + ograničenje sektora          & 0.68 & 0.21 & $-0.14$ & +0.26 & +0.36 & 0.97 \\
    + ogr.\ korelacijskih klastera & 0.71 & 0.20 & $-0.09$ & +0.31 & +0.26 & 0.86 \\
    \bottomrule
  \end{tabular}
\end{table}

Minimalna varijanca pokazuje jasnu pozitivnu izloženost faktorima RMW ($+0.33$) i CMA
($+0.34$) --- prepoznatljiv stil visoke profitabilnosti i niskoga ulaganja. Koncentracija
stila iznosi $0.98$, znatno iznad vrijednosti $0.83$ portfelja jednakih težina. Sektorsko
ograničenje smješteno je gotovo točno povrh neograničene minimalne varijance; korelacijsko
ograničenje koncentraciju umjereno snižava ($0.86$), no blok-bootstrap interval obuhvaća
nulu ($p\approx 0.99$), pa razlika nije razlučiva od šuma. Nagib je usto \emph{postojan}
kroz sve režime (slika~\ref{fig:concperyear}): oporavak nakon globalne financijske krize,
uzlet kvalitetnih dionica nakon 2014., šok pandemije COVID-19 te kamatni ciklus 2022.--2023.

\begin{figure}[H]
  \centering
  \includegraphics[width=0.92\linewidth]{04_style_concentration_per_year.png}
  \caption{Koncentracija stila po testnoj godini za pet izvornih portfelja. Nagib minimalne
  varijance postojan je kroz režime.}
  \label{fig:concperyear}
\end{figure}

Radi zornosti, slika~\ref{fig:weights} prikazuje kako svaki portfelj raspoređuje težinu po
sedam faktorskih klastera. Minimalna varijanca koncentrira približno 84\,\% težine u samo
dva klastera --- defenzivni klaster velikih dionica i klaster niskoga ulaganja --- čiji je
zajednički profil upravo RMW/CMA nagib koji je dijagnoza prepoznala.

\begin{figure}[H]
  \centering
  \includegraphics[width=0.92\linewidth]{06_cluster_weight_allocation.png}
  \caption{Raspodjela težina po sedam faktorskih klastera po portfelju (prozor 2024-12).}
  \label{fig:weights}
\end{figure}

\section{Čin II --- Ograničenje po klasteru i njegov neuspjeh}

Intuitivan ispravak glasi: budući da klasteri ocrtavaju stilska područja, ograničavanje
težine svakog klastera na 15\,\% trebalo bi prisiliti optimizator da se raširi cijelim
faktorskim prostorom i razrijedi nagib. Rezultat je, međutim, negativan
(tablica~\ref{tab:fail}).

\begin{table}[H]
  \centering
  \caption{Učinak ograničenja faktorskih klastera (puni uzorak).}
  \label{tab:fail}
  \small
  \begin{tabular}{lrrrrrr}
    \toprule
    Portfelj & $\beta_{\text{MKT}}$ & $\beta_{\text{SMB}}$ & $\beta_{\text{HML}}$ &
    $\beta_{\text{RMW}}$ & $\beta_{\text{CMA}}$ & konc.\ stila \\
    \midrule
    Minimalna varijanca           & 0.66 & 0.20 & $-0.12$ & +0.33 & +0.34 & 0.98 \\
    + ogr.\ faktorskih klastera   & 0.83 & 0.55 & $-0.19$ & $-0.15$ & +0.58 & \textbf{1.47} \\
    \bottomrule
  \end{tabular}
\end{table}

Ograničenje jest neutraliziralo ciljani RMW nagib ($+0.33\rightarrow-0.15$), no težinu je
premjestilo u klastere s jačom izloženošću faktorima SMB i CMA ($\beta_{\text{SMB}}:
0.20\rightarrow0.55$; $\beta_{\text{CMA}}: 0.34\rightarrow0.58$). Rezultat je portfelj s
\emph{većom} ukupnom koncentracijom stila ($1.47$) i za $5.8$ postotnih bodova višom
godišnjom volatilnošću (blok-bootstrap $p\approx 0.000$).

\paragraph{Zašto ispravak ne uspijeva.} Strukturni je razlog što su Wardovi klasteri u
prostoru faktora po konstrukciji zbijeni --- svaki je koherentna stilska skupina s
vlastitim nagibom. Ograničenje grupne težine nadzire \emph{koliko} težine pripada svakoj
skupini, ali ne i \emph{smjer} njezine faktorske izloženosti. Ukupna izloženost portfelja
težinski je prosjek tih nagiba, koji se ne moraju međusobno poništiti. Ključan je uvid
sljedeći: klasteri ispravno prepoznaju \emph{gdje} faktorska struktura leži (izvrsna su
dijagnostika), ali su pogrešan mehanizam za ograničavanje. Da bi se nagib doista smanjio,
ograničenje mora djelovati izravno na faktorske bete.

\section{Čin III --- Djelotvoran ispravak: izravna faktorska neutralnost}

Neuspjeh ograničenja po klasteru uputio je na pravo rješenje. Kvadratni program minimalne
varijance proširujemo linearnim ograničenjem koje djeluje upravo na geometriji koju je
klasteriranje otkrilo --- na faktorskim betama portfelja:
\[
  \big|\, w^{\!\top}\beta_f \,\big| \le \varepsilon,
  \qquad f \in \{\text{SMB},\ \text{HML},\ \text{RMW},\ \text{CMA}\}.
\]
Pritom je $\beta_f$ vektor faktorskih izloženosti dionica procijenjen na istom prozoru
procjene. $\varepsilon=0$ nameće strogu faktorsku neutralnost u trenutku procjene, dok
pozitivan $\varepsilon$ dopušta nadziran nagib. Ograničenje je linearno, funkcija cilja
ostaje nepromijenjena, a program konveksan.

\begin{table}[H]
  \centering
  \caption{Atribucija i ostvarena volatilnost faktorski neutralne obitelji (puni uzorak).}
  \label{tab:fix}
  \small
  \begin{tabular}{lrrrrrrr}
    \toprule
    Portfelj & $\beta_{\text{MKT}}$ & $\beta_{\text{SMB}}$ & $\beta_{\text{HML}}$ &
    $\beta_{\text{RMW}}$ & $\beta_{\text{CMA}}$ & konc.\ stila & god.\ vol. \\
    \midrule
    Minimalna varijanca           & 0.66 & +0.20 & $-0.12$ & +0.33 & +0.34 & 0.98 & 12.0\,\% \\
    faktorski neutralan, $\varepsilon{=}0.00$ & 0.69 & +0.18 & $-0.03$ & +0.26 & +0.20 & \textbf{0.67} & 12.1\,\% \\
    faktorski neutralan, $\varepsilon{=}0.05$ & 0.68 & +0.19 & $-0.06$ & +0.29 & +0.22 & 0.76 & 12.0\,\% \\
    faktorski neutralan, $\varepsilon{=}0.10$ & 0.67 & +0.19 & $-0.07$ & +0.31 & +0.24 & 0.81 & 12.0\,\% \\
    faktorski neutralan, $\varepsilon{=}0.15$ & 0.67 & +0.21 & $-0.11$ & +0.30 & +0.31 & 0.93 & 12.0\,\% \\
    \bottomrule
  \end{tabular}
\end{table}

Najstroža postavka $\varepsilon=0$ smanjuje ostvarenu koncentraciju stila za 31\,\%
($0.98\rightarrow0.67$) uz zanemariv trošak u volatilnosti. Blok-bootstrap interval razlike
koncentracije stila u odnosu na minimalnu varijancu u cijelosti je ispod nule ($-0.31$;
95\,\% CI $[-0.47,\,-0.09]$; $p=0.002$), dok interval razlike volatilnosti obuhvaća nulu
($p=0.59$). Upravo je to kompromis koji je dijagnoza tražila (slika~\ref{fig:bootstrap}).

\begin{table}[H]
  \centering
  \caption{Blok-bootstrap intervali razlika u odnosu na \texttt{min\_var} (1000 ponovnih uzoraka).}
  \label{tab:boot}
  \small
  \begin{tabular}{lrrr}
    \toprule
    Usporedba & $\Delta$ konc.\ stila & 95\,\% CI & $p$ \\
    \midrule
    faktorski neutralan $\varepsilon{=}0.00$ & $\mathbf{-0.308}$ & $[-0.47,\,-0.09]$ & \textbf{0.002} \\
    faktorski neutralan $\varepsilon{=}0.05$ & $-0.225$ & $[-0.39,\,-0.06]$ & 0.004 \\
    faktorski neutralan $\varepsilon{=}0.10$ & $-0.171$ & $[-0.32,\,-0.03]$ & 0.006 \\
    faktorski neutralan $\varepsilon{=}0.15$ & $-0.054$ & $[-0.10,\,-0.01]$ & 0.020 \\
    \bottomrule
  \end{tabular}
\end{table}

\begin{figure}[H]
  \centering
  \includegraphics[width=0.62\linewidth]{05_bootstrap_factor_neutral_e0.png}
  \caption{Bootstrap-distribucija razlike koncentracije stila (faktorski neutralan
  $\varepsilon{=}0$ $-$ \texttt{min\_var}). Interval leži u cijelosti ispod nule.}
  \label{fig:bootstrap}
\end{figure}

Granica koncentracije stila naspram ostvarene volatilnosti (slika~\ref{fig:frontier})
smješta čistu faktorski neutralnu obitelj na efikasnu granicu: ona smanjuje nagib bez
porasta volatilnosti, dok hibridni portfelji (klaster $+$ neutralnost) i faktorsko-klasterski
portfelj to plaćaju višom volatilnošću.

\begin{figure}[H]
  \centering
  \includegraphics[width=0.72\linewidth]{05_frontier_style_vs_vol.png}
  \caption{Koncentracija stila naspram ostvarene godišnje volatilnosti za sve obitelji
  portfelja. Faktorski neutralna obitelj zauzima efikasnu granicu.}
  \label{fig:frontier}
\end{figure}

\section{Pokrivenost metoda kolegija}

Dva dodatka čvršće povezuju projekt s metodama kolegija i podvrgavaju njegove odluke
dodatnoj provjeri.

\paragraph{Algoritmi klasteriranja i odabir $K$ (predavanja 9b/10).} Prostor faktora
ponovno smo klasterirali particijskim metodama (K-means, k-medoid), vjerojatnosnom
metodom (mješavina Gaussovih distribucija, EM) te metodom utemeljenom na gustoći (DBSCAN),
uz odabir $K$ pomoću siluete, Calinski--Harabaszova i Davies--Bouldinova indeksa te Gap
statistike (slika~\ref{fig:kselect}). Zaključci: (i) struktura je otporna na izbor
algoritma --- K-means daje u osnovi iste skupine kao Wardova metoda (ARI $\approx 0.45$);
(ii) sva četiri indeksa interne validacije upućuju na \emph{statistički} prirodan $K$
između 3 i 5, pa je $K=7$ ponajprije uvjet uloživosti, a ne optimum siluete --- što dijelom
objašnjava zašto je ispravak ograničenjem po klasteru imao tako malo strukture za
iskoristiti; (iii) DBSCAN ne pronalazi prirodne praznine u gustoći (ARI $\approx 0$), čime
se potvrđuje da su metode sa zadanim $K$ ovdje pravi alat.

\begin{figure}[H]
  \centering
  \includegraphics[width=0.95\linewidth]{07_kselection_indices.png}
  \caption{Indeksi za odabir broja klastera i Gap statistika za Wardovu metodu i K-means.
  Prirodan $K$ nalazi se u rasponu 3--5.}
  \label{fig:kselect}
\end{figure}

\paragraph{Nadzirano proširenje --- predviđanje pada (predavanja 3--4).} Za svaku dionicu i
testno razdoblje označavamo „visok pad” (gori od medijana u presjeku te godine) te ga
predviđamo iz značajki prozora procjene --- stablima odlučivanja (entropija / Gini, dubina
3) i ansamblima (slučajna šuma, gradijentni boosting) --- uz unakrsnu validaciju grupiranu
po testnoj godini (GroupKFold).

\begin{table}[H]
  \centering
  \caption{Unakrsna validacija predviđanja pada (najbolji model po skupu značajki).}
  \label{tab:cv}
  \small
  \begin{tabular}{lrr}
    \toprule
    Skup značajki & najbolji ROC-AUC & najbolja točnost \\
    \midrule
    samo faktorske bete       & 0.69 & 0.64 \\
    + GICS sektor             & 0.69 & 0.63 \\
    + oznaka faktorskog klastera & 0.68 & 0.64 \\
    \bottomrule
  \end{tabular}
\end{table}

Pad je umjereno predvidljiv (AUC $\approx 0.69$ naspram 0.5), no signalom dominiraju
rezidualna volatilnost (idiosinkratski rizik) i tržišna beta, a ne stilski faktori.
Dodavanje GICS sektora ili oznake klastera ne donosi poboljšanje u odnosu na same bete:
klaster je izvrstan \emph{opisni} objekt, ali ne i dodatna \emph{prediktivna} značajka.
Taj rezultat iznosimo onakvim kakav jest --- opisna struktura ne služi ujedno i kao
prediktor.

\section{Zaključak}

\textbf{Otkriće.} Portfelji minimalne varijance na indeksu Russell 1000 nose skriven i
postojan nagib prema faktorima profitabilnosti (RMW) i ulaganja (CMA) --- koncentraciju
stila skrivenu sektorskoj i korelacijskoj diverzifikaciji, ali jasno ocrtanu klasteriranjem
u prostoru faktora.

\textbf{Ispravak.} Dodavanje linearnih ograničenja $|w^{\!\top}\beta_f|\le\varepsilon$
izravno u kvadratni program smanjuje ostvarenu koncentraciju stila za 31\,\% pri
$\varepsilon=0$, bez vidljiva troška u volatilnosti ($p\approx 0.002$).

\textbf{Pouka.} Struktura faktorskih klastera bila je ispravna dijagnostika --- učinila je
skriveni nagib vidljivim --- ali pogrešan mehanizam ispravka: ograničavanje grupne težine
ne cilja nijednu pojedinačnu betu. Djelotvoran ispravak postavlja ograničenje točno ondje
gdje je klasteriranje pokazalo problem: na same faktorske bete.

\paragraph{Glavna ograničenja.} (i) Pristranost preživljavanja --- univerzum je trenutačni
sastav indeksa Russell 1000, a ne indeks u stvarnoj povijesnoj točki. (ii) Bete su
procijenjene iz 60 mjesečnih opažanja, dovoljno malo da unesu šum u procjenu. (iii)
Transakcijski troškovi i obrtaj zanemareni su. (iv) Predviđanje prinosa nije cilj; brojevi
kumulativnoga rasta nuspojava su, a ne rezultat. (v) Korelacijsko-klasterski portfelj i
hibridi pri $\varepsilon\in\{0,0.05\}$ agregirani su preko manjega broja testnih mjeseci
(168/180 naspram 192) zbog neizvedivosti u najranijim prozorima.

\vspace{1em}
\hrule
\vspace{0.6em}
{\small Programski kod, podatci i sve slike: repozitorij projekta (\texttt{src/},
\texttt{notebooks/} 01--08, \texttt{outputs/}). Detaljan narativni izvještaj:
\texttt{reports/final\_report.md}.}

\end{document}
